\documentclass[a4paper,10pt]{report}\input{../0.Préambule/Préambule_cours_prof}\begin{document}

\newpage
\section*{Grandeurs composées}
\addcontentsline{toc}{section}{Grandeurs composées}

\vspace{-4mm}
\begin{exer1}
Un automobiliste parcourt $72$km en $54$min. Quelle est sa vitesse moyenne en km/h ?
\end{exer1}

\begin{exer1}
Un cycliste roule $2$h$35$min à la vitesse moyenne de $27$km/h. Quelle distance a-t-il parcourue ?
\end{exer1}

\begin{exer1}
Dans une commune, $43$m$^3$ d'eau sont facturés $107,07$\euro. Quel est le prix de l'eau en \euro /m$^3$ ?
\end{exer1}

\begin{exer1}
A Saint-Malo, Erwan a prélevé $1,25$L d'eau de mer qu'il a fait évaporer et il a obtenu $43$g de sel. Quelle est la concentration en sel de la Manche à Saint-Malo, en g/L ?
\end{exer1}

\begin{exer1}
Le département de l'Isère compte 1,264 millions d'habitants répartis sur 7 431 km$^2$. Quelle est la densité de l'Isère en habitants/km$^2$ ?
\end{exer1}

\begin{exer1}
Combien de temps faut-il pour parcourir $800$m à la vitesse moyenne de $40$km/h ? Exprime le résultat en minutes secondes.
\end{exer1}

\begin{exer1}
Jules court $100$m en $12,5$secondes. Calcule sa vitesse en m/s puis en km/h.
\end{exer1}

\begin{exer1}
Pour monter un mur, Stéphane a utilisé $300$kg de béton dont la masse volumique est $2,5$ tonnes par mètre cube. Quel volume de béton a-t-il utilisé ?
\end{exer1}

\begin{exer1}
Une imprimante a imprimé $36$ pages en $1$min $48$s. Quelle est la vitesse d'impression en page/min ?
\end{exer1}

\begin{exer1}
Une fraise de dentiste tourne de $\np{1 500}$ tours par minute. Exprime cette vitesse en tours/seconde.
\end{exer1}

\begin{exer1}
Une perfusion délivre $1,5$L de solution pendant $24$h. Quel est le débit de la perfusion en mL/min ?
\end{exer1}

\begin{exer1}
Le débit d'un robinet est de $210$L/h. Quelle quantité d'eau coule en le laissant ouvert $3$min ?
\end{exer1}

\begin{exer1}
Noémie a fait le plein et a mis $60$L d'essence dans son réservoir. La masse volumique de l'essence est de 750 g/m$^3$. De quelle masse son véhicule a-t-il augmenté ?
\end{exer1}

\begin{exer1}
A $12$km/h, en combien de temps parcourt-on $1$km ? Donne le résultat dans une unité adaptée.
\end{exer1}

\begin{exer1}
Le train Eurostar permet de relier la France et l'Angleterre en passant dans le tunnel sous la Manche long de $50,5$km. Il parcourt le tunnel en $29$min. Quelle est sa vitesse moyenne dans le tunnel en km/h ?
\end{exer1}

\begin{exer1}
La vitesse de la lumière est $\np{300 000}$km/s. La lumière met environ $8$minutes et $30$secondes pour nous parvenir du Soleil. Calcule la distance nous séparant du Soleil.
\end{exer1}

\begin{exer1}
Un cycliste part de chez lui et roule à la vitesse moyenne de $15$km/h pendant $30$km pour monter un col. Il fait ensuite demi-tour pour rentrer chez lui à la vitesse moyenne de 30 km/h. Calcule sa vitesse moyenne sur la totalité de son trajet aller-retour.
\end{exer1}

\begin{exer1}
La paille est vendue au prix de gros de $40$\euro~ la tonne. Un mètre cube de paille pèse $120$kg. Les bottes de paille sont des pavés droits de dimensions $80$cm $\times$ $40$cm $\times$ $30$cm. Calcule le prix d'une botte de paille.
\end{exer1}

\begin{exer1}
Combien de temps un automobiliste va gagner en roulant à la vitesse moyenne de $110$km/h au lieu de $80$km/h sur un trajet de $35$km ?
\end{exer1}

\begin{exer1}
On remplit une citerne de $175$L avec une pompe dont le débit est $2,1$m$^3$/h. Combien faut-il de temps pour remplir cette citerne ?
\end{exer1}

\begin{exer1}
On considère la fenêtre de téléchargement ci-contre.

\vspace{-0.7cm}
\begin{flushright}
\includegraphics[scale=0.37]{../40.Image/12.Exo-Grandeurs_composées_1}
\end{flushright}

\vspace{-0.7cm}
\noindent Si la vitesse de téléchargement reste constante, faudra-t-il plus d'une \\minute et vingt-cinq secondes pour que le téléchargement se termine ?
\end{exer1}

\begin{exer1}
Une tige d'acier a la forme d'un cylindre qui mesure $1,5$m de long et a pour rayon de base $4$cm. Sachant que l'acier a une masse volumique de $7,85$g/cm$^3$, calcule la masse de cette tige.
\end{exer1}

\end{document}